\chapter{Análisis Forense de Cuentas Comprometidas}

La seguridad de un sistema de cómputo multiusuario es tan sólida como el eslabón más débil de su cadena de credenciales. Cuando un atacante remoto logra adivinar o interceptar la contraseña de una cuenta legítima, no requiere explotar complejas vulnerabilidades de corrupción de memoria (\textit{buffer overflows}) para ingresar al servidor: entra por la puerta principal amparado por los privilegios asignados al usuario.

En este capítulo se expone la inspección pericial del perfil del usuario \texttt{proa} (UID 523), la cuenta utilizada como vector de penetración inicial y depósito logístico del arsenal de herramientas del atacante.

\section{El Depósito Clandestino en el Perfil de proa}

Tal como se constató en las bitácoras de transferencia FTP analizadas en el Capítulo 9, la cuenta \texttt{proa} fue utilizada para introducir los paquetes de código fuente malicioso en el servidor. Al inspeccionar el contenido del directorio personal \texttt{/home/proa}, el equipo pericial halló un inventario extraordinario de artefactos ofensivos, resumido en el Listado~\ref{lst:proa_home_curated}.

\begin{lstlisting}[style=terminal,caption={Extracto representativo del inventario de artefactos hallados en /home/proa.},label={lst:proa_home_curated}]
drwxr-xr-x 10 523      523         4096 Apr 16 22:16 ...
drwxr-xr-x  3 523      523         4096 Apr 16 21:33 .king
drwx------  4 523      523         4096 Apr 15 22:20 emech-2.8.3
drwxr-xr-x  8 523      523         4096 Apr 12 05:33 psybnc
-rw-------  1 523      523          404 Apr  9 20:19 .bash_history
-rw-r--r--  1 523      523       200419 Apr 20 00:00 EriN[ca].seen
-rw-r--r--  1 523      523       184395 Apr 20 00:00 victima2.seen
-rw-r--r--  1 523      523          143 Apr 20 00:00 bot.users
-rwxr-xr-x  1 523      523       466184 Apr 15 23:39 in.ftpd
-rw-r-----  1 523      523        37210 Apr 16 10:01 passwd
-rw-r--r--  1 523      523       970089 Oct 28  2000 psyBNC2.2.1-linux-i86-static.tar.gz
-rw-r--r--  1 523      523      1546240 Feb 13 10:13 psyBNC2.3.1.tar
-rw-r--r--  1 523      523           77 Mar  9 23:40 psybnc.conf
\end{lstlisting}

\subsection{Categorización Pericial de los Objetos Descubiertos}

El Cuadro~\ref{tab:analisis_proa} clasifica las herramientas localizadas en el perfil de acuerdo con su finalidad técnica.

\begin{table}[htbp]
\centering
\small
\begin{tabular}{lp{10cm}}
\toprule
\textbf{Categoría} & \textbf{Artefactos Hallados y Función Pericial} \\
\midrule
\textbf{Directorios Ocultos} & \texttt{...} (directorio de tres puntos) y \texttt{.king}. Estructuras camufladas para almacenar ejecutables y scripts de inicio automático. \\
\addlinespace
\textbf{Bots de Red IRC} & \texttt{emech-2.8.3} (\textit{EnergyMech}), \texttt{bot.users}, \texttt{EriN[ca].seen}, \texttt{victima2.seen}. Distribución completa del bot EnergyMech y bases de datos de seguimiento de nicks en canales IRC. \\
\addlinespace
\textbf{Bouncers (Proxies IRC)} & \texttt{psybnc}, versiones 2.2.1 y 2.3.1 en código fuente y paquetes estáticos precompilados, junto con el archivo de configuración activo \texttt{psybnc.conf}. \\
\addlinespace
\textbf{Servicios Clandestinos} & \texttt{in.ftpd}. Demonio de FTP independiente compilado para ejecutarse en puertos no estándar, permitiendo la carga y descarga de archivos eludiendo el servicio FTP oficial del sistema. \\
\addlinespace
\textbf{Sustracción de Credenciales} & \texttt{passwd} (archivo de 37.2~KB). Copia local del archivo de cuentas del sistema, sustraída para la ejecución de ataques de fuerza bruta locales (\textit{offline cracking}) con herramientas como \textit{John the Ripper}. \\
\bottomrule
\end{tabular}
\caption{Clasificación funcional de las herramientas descubiertas en la cuenta proa.}
\label{tab:analisis_proa}
\end{table}

\section{Técnicas de Ocultamiento: El Directorio «Tres Puntos» (\texttt{...})}

Uno de los mecanismos de camuflaje más antiguos y extendidos en la cultura del \textit{cracking} sobre sistemas Unix fue el uso de directorios denominados literalmente \texttt{...} (tres puntos consecutivos).

En los sistemas de archivos compatibles con POSIX, el comando \texttt{ls -a} muestra obligatoriamente las entradas especiales \texttt{.} (directorio actual) y \texttt{..} (directorio padre). Al crear un subdirectorio nombrado con tres puntos (\texttt{mkdir ...}):
\begin{itemize}
    \item En listados apresurados, el ojo del administrador tiende a confundir la entrada \texttt{...} con los metadirectorios estándar del sistema.
    \item Los scripts de limpieza basados en expresiones regulares simples o que filtran directorios que comienzan con punto pueden omitir o manejar erróneamente esta ruta.
\end{itemize}

La inspección del subdirectorio \texttt{/home/proa/...} confirmó que operaba como el centro neurálgico del bot \textit{EnergyMech}, conteniendo sus ejecutables principales, registros de conexión y scripts en shell para reiniciar el proceso de forma persistente cada vez que el sistema se reiniciaba.

\section{Disección del Bot EnergyMech (emech-2.8.3)}

\textit{EnergyMech} representó una de las evoluciones más avanzadas en la programación de bots para IRC en C durante las décadas de 1990 y 2000. A diferencia de los clientes tradicionales, incorporaba soporte nativo para:
\begin{itemize}
    \item \textbf{Gestión jerárquica de usuarios (\texttt{bot.users} y \texttt{mech.levels}):} Permitía asignar niveles de mando (maestro, administrador, operador) a los usuarios remotos que enviaran comandos autenticados a través de mensajes privados en IRC.
    \item \textbf{Bases de datos de presencia (\texttt{.seen}):} Los archivos \texttt{EriN[ca].seen} y \texttt{victima2.seen} almacenaban el registro de todas las personas que entraban y salían de los canales de chat monitoreados por el bot, guardando su máscara de host, fecha y último mensaje emitido.
    \item \textbf{Capacidades ofensivas integradas:} Módulos de ataque para saturación de canales mediante envío masivo de mensajes (\textit{flooding}) y desestabilización de clientes mediante cadenas maliciosas.
\end{itemize}

\section{Exfiltración Local de Credenciales: El Archivo passwd}

El hallazgo de un archivo de 37,210 bytes titulado \texttt{/home/proa/passwd} aportó una prueba crítica sobre las intenciones del intruso. Al cotejar su contenido con el archivo original \texttt{/etc/passwd} del servidor, se comprobó que se trataba de una réplica íntegra de la base de usuarios de la universidad.

En los sistemas Unix configurados con contraseñas ocultas (\textit{shadow passwords}), las cadenas hash cifradas residen exclusivamente en \texttt{/etc/shadow}, un archivo con permisos de lectura reservados al superusuario (\texttt{0400} de \texttt{root}). Sin embargo, la copia hallada en el directorio de \texttt{proa} contenía las entradas completas de nombres de usuario, directorios personales, intérpretes de comandos asignados y números de identificación (UID/GID). 

Esta información suministró al atacante el mapa completo de usuarios del centro de cómputo, facilitándole la selección de cuentas candidatas para ejecutar ataques de diccionario dirigidos.

\section{Vectores de Débil Autenticación y Factores Sistémicos}

La investigación forense de las cuentas afectadas permitió identificar las causas estructurales que hicieron posible el compromiso:
\begin{enumerate}
    \item \textbf{Contraseñas triviales:} La cuenta \texttt{proa} utilizaba una contraseña idéntica al propio nombre de usuario o constituida por palabras del diccionario español sin números ni caracteres especiales, vulnerable a ser descubierta en cuestión de segundos mediante escaneos de fuerza bruta automatizados sobre el puerto 21 (FTP).
    \item \textbf{Ausencia de políticas de bloqueo de cuentas:} En la configuración de red de 2004, los servicios no limitaban el número de intentos fallidos consecutivos ni existían sistemas de baneo automático de direcciones IP como \textit{Fail2ban}.
    \item \textbf{Protocolos en texto claro:} La provisión de servicios sin cifrado (FTP en puerto 21 y Telnet en puerto 23) facilitó que las credenciales fueran interceptadas en tránsito cada vez que los usuarios legítimos iniciaban sesión desde terminales externas.
\end{enumerate}

\section{Conclusiones del Análisis de Cuentas}

El examen forense del perfil \texttt{proa} demostró que:
\begin{itemize}
    \item La cuenta no fue creada por el intruso, sino que pertenecía a un usuario legítimo cuya clave fue vulnerada desde el exterior.
    \item La cuenta funcionó como la base logística y depósito de almacenamiento (\textit{drop zone}) donde se compilaron y ejecutaron los servicios de IRC y bouncers antes de que el ataque secundario escalara a nivel de superusuario.
\end{itemize}
